\begin{theorem}
Let \(k\) be an algebraically closed field with the trivial absolute value,
let \(n\geq2\), and put
\[
             X=\mathbf P^n_k,
             \qquad L=\mathcal O_{\mathbf P^n}(1),
             \qquad \omega=c_1(L).
\]
Fix a linear line \(C\simeq\mathbf P^1_k\) in \(X\).  For \(n\geq3\),
let \(v_C=\operatorname{ord}_{E_C}\), where \(E_C\) is the exceptional
divisor of the blowup of \(X\) along \(C\); for \(n=2\), set
\(v_C=\operatorname{ord}_C\).  Choose pairwise distinct points
\(p_j\in C(k)\), and let
\(v_j=\operatorname{ord}_{E_j}\), where \(E_j\) is the exceptional divisor
of the blowup of \(X\) at \(p_j\).  Then
\[
              v_j\longrightarrow v_C
              \quad\text{in the weak topology of }X^{\mathrm{lin}},
\]
and
\[
              T_\omega(v_j)=T_\omega(v_C)=1
              \qquad\text{for every }j,
\]
whereas
\[
              d_{\infty,\omega}(v_j,v_C)\geq1
              \qquad\text{for every }j.
\]
Consequently, weak convergence together with convergence of the
Boucksom--Jonsson exhaustion does not imply strong convergence on
\(X^{\mathrm{lin}}\).  This failure occurs on \(\mathbf P^n_k\) for every
\(n\geq2\), in arbitrary characteristic.
\end{theorem}
